\chapter{Related Work}\label{chap:relatedWork}
Our work lies at the intersection of two lines of research: 
cost analysis and productivity of infinite programs. 
In this chapter, we discuss related work from both areas, 
focusing on static approaches, particularly type-based techniques, 
and discuss their connections to our work.

\section{Cost Analysis}
\label{sec:rw-cost}
Techniques for (type-based) static cost analysis are commonly divided into two styles
\citep{rajaniUnifyingTypetheoryHigherorder2021a}. In \emph{effect-based}
systems, cost can be considered as an extra output of a computation, as described in
Chapter~\ref{chap:background}. In \emph{coeffect-based} systems, cost is
instead treated as an additional requirement of programs, such
as the number of times each input can be used. Specifically, affine types used in $\lambda$-amor 
and our $\lambda^g_{\mathsf{amor}}$ are coeffects.

\subsubsection{Effect-Based Cost Analysis}
Automatic Amortised Resource Analysis (AARA) is one of the most extensively developed approaches to type-based resource analysis.
First presented by \citet{hofmannJ03}, it employs a type system based on the potential method to automatically infer linear bounds on the resource consumption of first-order functional programs. 
Subsequent work extended AARA to support univariate and multivariate polynomial potential \citep{hoffmannH10, hoffmannMultivariateAmortizedResource2012a}, 
in which resource bounds can depend on multiple program inputs. 
To prevent the unsound duplication of potential, the higher-order extension of AARA by \citet{JostHLH10} restricts function closures from capturing potential, 
allowing functions to be applied arbitrarily many times. However, this restriction still limits the analysis of higher-order functions whose costs depend on the data they carry.

One approach to address this limitation is to use affine types, as in $\lambda$-amor \citep{rajaniUnifyingTypetheoryHigherorder2021a} and our work,
to control the duplication of potential rather than prohibiting capturing potential.
By contrast, \citet{xuW26} proposed a variant of AARA that separates potential from types to traditional potential functions via dependent types, 
and hence does not rely on affine types to analyse higher-order programs.

Another line of work on effect-based cost analysis is recurrence extraction \citep{dannerDenotationalCostSemantics2015,kavvosRecurrenceExtractionFunctional2020,cutlerDenotationalRecurrenceExtraction2020}, 
which provides an alternative to deriving cost bounds directly through type systems. 
These approaches translate programs into cost recurrences that can subsequently be solved to obtain their complexity bounds. 
The translation uses a writer monad, which makes cost an explicit output of computations. 
This is closely related to our use of a graded monad,  
although their costs are represented explicitly in the translated programs rather than tracked through the grades contained by types.

\subsubsection{Coeffect-Based Cost Analysis}
An early precursor to coeffect-based resource analysis is Bounded Linear Logic (BLL)~\citep{girardSS92}.
BLL derives cost by tracking variable usage through affine types. Unlike the standard affine type system adopted in $\lambda^g_{\mathsf{amor}}$, where variables can be used at most once unless assigned an exponential type 
$!\tau$, which permits unrestricted usage, BLL uses a graded exponential type $!_I\tau$, where the grade $I$ tracks an upper bound on variable usage. 
By explicitly bounding variable duplication, it provides a logical foundation for polynomial-time cost analysis.

$d\ell\mathrm{PCF}$ \citep{dallagoLinearDependentTypes2011a} is a type system based on BLL for analysing the cost of 
call-by-name PCF programs, including recursion and higher-order functions \citep{Plotkin77}.
\cite{gaboardiCombiningEffectsCoeffects2016} develop a unified type system for effects and coeffects, using graded modal types to capture both of them.
This is related to the approach of $\lambda$-amor and our work, which also employ both graded effect types (monads) and coeffect types (affine types and potential).
This approach is further explored in a practical functional programming language, Granule \citep{orchardQuantitativeProgramReasoning2019a}. Although the above systems can express both 
effects and coeffects, their main focus is tracking resource usage and ownership through the latter.

The full $\lambda$-amor calculus can also embed $d\ell\mathrm{PCF}$ through its graded exponential type $!_{I}\tau$. 
In contrast, our calculi only include the ordinary exponential type $!\tau$, 
as we focus on effect-based cost analysis rather than coeffects, while $\lambda$-amor provides a unified framework for both approaches.

\section{Type-Based Productivity}
\label{sec:rw-productivity}

\subsection{Sized Types}
Besides guarded recursion, another well-studied type-based approach to ensuring productivity is \emph{sized types},
where recursive types are annotated with size indices \citep{hughesPS96}. 
For example, a sized stream type $\mathsf{Stream}^{k}(\tau)$ guarantees that at least the first $k$ elements can be observed.
This notion of observation depth is closely related to the evaluation fuel $n$ in our operational semantics, 
which specifies the number of guarded computation layers that can be explored. However, sized types express such guarantees statically through the type system,
whereas our evaluation fuel controls finite observations operationally.

\subsection{Guarded Recursion}
\label{sec:guarded-related-work}
Several extensions of guarded recursion make the notion of logical
time explicit through clock-indexed later modalities
\citep{atkeyProductiveCoprogrammingGuarded2013a,mogelberg14,
bahrClocksAreTicking2017}.
These systems support multiple clocks, allowing different
computations to progress independently in logical time.
For example, a pair of type
$\triangleright^\alpha \tau_1 \times \triangleright^\beta \tau_2$
contains two components whose availability is controlled by
different clocks $\alpha$ and $\beta$.
In contrast, we follow the guarded $\lambda$-calculus
\citep{cloustonGuardedLambdaCalculusProgramming2017a},
which just uses a single global clock.
This choice is particularly suitable for our work, where
costs may depend on logical time, as it ensures that
all costs vary uniformly with respect to the
same clock.

\cite{guattoGeneralizedModalityRecursion2018d} generalises the later modality to a warping modality $_p\tau$, where $p$ is a time warp function that reindexes the logical time, 
determining how much information is available at each stage of computation. 
The ordinary later modality is recovered by a one-step delay, written $_{-1}\tau$. 
More generally, time warps can express different rates of information production, 
such as streams that produce a new element every two time steps. 
This provides more fine-grained control over temporal dependencies in guarded recursion. 
A similar notion of temporal reindexing appears at the level of cost indices in our calculi: $\mathsf{S}^{s}(I)$ interprets a time-dependent cost index $I$ at $s$ steps into the future.

The interaction between guarded recursion and computational effects has also been studied. 
\cite{mogelbergWhatMonadsCan2025} investigate distributive laws between the delay monad,
which uses guarded recursion to represent computations that may either terminate with a value or continue beneath a later modality,
and other effect monads, including the writer monad that can model costs.
A similar interaction arises in our calculi through the $\mathsf{prepay}$ operator, 
which transforms $\triangleright\tmonad{I}{\tau}$ into $\tmonad{S(I)}{\triangleright\tau}$. 
This corresponds to a distributive law between the later modality and the graded monad, although we only support one direction: there is no corresponding operator that transforms $\tmonad{I}{\triangleright\tau}$ into $\triangleright\tmonad{J}{\tau}$.

\section{Logical Approaches to Cost Analysis}
\label{sec:rw-logics}
Cost can also be verified by stating its properties as propositions and proving
them using logics, rather than checking them automatically with a type system.
\cite{niuCostawareLogicalFramework2022a} introduce CALF, 
a cost-aware logical framework based on dependent type theory that treats cost as a computational effect. 
CALF allows both the functional behaviour and resource usage of programs to be specified and verified within the same language. 
In particular, cost bounds are proved through internal logical reasoning, and the framework supports potential-based amortised analysis.
Based on CALF, \cite{grodinAmortizedAnalysisCoinduction2023a} propose a coinductive formulation of amortised analysis, 
which reasons about potentially infinite behaviour of data structures under continuous operations, instead of analysing finite sequences of operations. 
This work is related to $\lambda^g_{\mathsf{amor}}$ in reasoning about infinite computations and data structures through amortised cost analysis; the difference is that we mainly focus on type-based analysis.

In addition to logical frameworks, a more common alternative is program logic.
In particular, separation logic \citep{Reynolds02} has been widely adopted for its support of reasoning about program resources, such as mutable memory.
\cite{Atkey10} extends separation logic with resource assertions to support amortised analysis of imperative programs. There has also been further work on potential-based amortised analysis with separation logic \citep{ChargueraudP19, MevelJP19}.
Among these approaches, the potential available before and after an operation can be represented as ghost state, which is similar to 
the potential used by the type systems and logical relations of $\lambda$-amor and $\lambda^g_{\mathsf{amor}}$.

\section{Cost Analysis of Lazy Programs}
\label{sec:rw-lazy}
Lazy evaluation, also known as call-by-need evaluation, suspends a computation until
its result is demanded and shares the result after evaluation
\citep{launchburyNaturalSemanticsLazy1993a}. This makes it particularly suitable for working with infinite data structures. 
The banker's method of amortised analysis uses credits accumulated during earlier operations to pay for future costs. \cite{okasak} adapts this method to lazy data structures by using \emph{debits} to represent the unpaid costs of suspended computations. These debits are gradually discharged by other operations, ensuring that the cost of a suspended computation is accounted for before it is evaluated.
Moreover, unlike potential-based approaches, such as $\lambda$-amor, amortised analysis in this setting does not require affine types, since shared computations are evaluated at most once

Following Okasaki's approach, \cite{Danielsson08} develops a type-based framework for analysing lazy functional programs using a graded monad to track amortised costs. In particular, it introduces a primitive operation $\mathsf{pay}$, which, written in our notation, has type
$\tmonad{I_1+I_2}{\tau} \to \tmonad{I_1}{(\tmonad{I_2}{\tau})}$.
This operator allows part of the deferred cost of a suspended computation to be paid in advance, reducing the remaining cost. 
This may seem to be similar to our $\mathsf{prepay}$ operator, which transforms $\triangleright\tmonad{I}{\tau}$ into $\tmonad{S(I)}{\triangleright\tau}$. However, $\mathsf{pay}$ redistributes amortised costs between nested lazy computations to implement Okasaki's approach, while our $\mathsf{prepay}$ moves a cost 
across the later modality, accounting for its dependence on logical time associated with the modality.

Based on AARA, \cite{jostTypeBasedCostAnalysis2017a} present a cost analysis type system
for lazy functional programs, where lazy computations are given
thunk types annotated with their costs. The analysis accounts for the
sharing of thunks and derives precise upper bounds on the evaluation cost of lazy
programs, including potentially non-terminating programs such as processing infinite lists.

In comparison, our work uses call-by-value evaluation, where the graded monad suspends a
computation similar to the thunk type used by \cite{jostTypeBasedCostAnalysis2017a}, but the result after forcing a monad is not shared.
We discuss a possible extension towards call-by-need evaluation in
Chapter~\ref{chap:conclusion}.

\section{Modal Types for Time and Resource Analysis}
\label{sec:rw-temporal}
There are other works that use modal types to describe the
relationship between time, resources and costs in programs.

Functional reactive programming (FRP) describes programs that react to events and produce values that change over time continuously \citep{elliottH97}. 
Modal types and guarded recursion have been used in FRP to ensure productivity and prevent undesirable behaviours such as space leaks, where programs retain unnecessary data in memory \citep{krishnaswamiBH12,bahrSimplyRaTTFitchstyle2019a}. 
While both approaches employ guarded recursion, these works focus on controlling resource usage in reactive programs, whereas our work aims to derive quantitative cost bounds for those programs.

In concurrent programming, session types provide a type-based framework for specifying communication protocols \citep{HondaVK98}. 
\cite{dasParallelComplexityAnalysis2018} extend session types with temporal modalities to specify when messages can be exchanged, 
allowing the type system to derive bounds on the parallel execution time of programs. 
In their approach, computational costs are measured directly in units of logical time, such as latency.
In contrast, our calculi treat logical time and computational cost separately. Despite the dependency between time and costs, costs are not directly represented by time in our work.

